\chapter{HTML}

\section{Dentro un server HTTP}

Nel capitolo precedente abbiamo visto che il Web è un sistema di documenti
ipertestuali collegati fra loro, e che i client usano HTTP per recuperarli.
Prima di scrivere il nostro primo documento, chiariamo che cosa c'è dall'altra
parte del filo.

Un \term{server HTTP} è, molto semplicemente, un \textbf{programma}. Le sue
responsabilità sono tre:

\begin{itemize}
  \item si mette in \textbf{ascolto} di richieste HTTP su una certa porta (per
        esempio la 80);
  \item può gestire \textbf{più host} contemporaneamente: è esattamente per
        questo che nelle richieste HTTP esiste l'header \code{Host}, che dice al
        server \textit{quale} dei siti ospitati stiamo cercando;
  \item risponde alle connessioni servendo file prelevati dalla propria
        \term{document root} nel filesystem.
\end{itemize}

\begin{figure}[H]
  \centering
  \begin{tikzpicture}[font=\Lato\small,
      host/.style={rounded corners=4pt, draw=primary!55!black, line width=0.8pt,
                   fill=primary!5!white},
      f/.style={wtnodeplain, font=\FiraCode\scriptsize, minimum height=0.5cm},
      d/.style={f, draw=orange-gradient-6, fill=orange-gradient-1!40!white}]

    \node[font=\Lato\bfseries, text=primary!75!black] (t) at (0,3.3) {SERVER HTTP};

    % --- host 1 ---
    \node[font=\Lato\scriptsize, anchor=west, text=primary!45!black] (h1t)
      at (-5.1,2.45) {Host: \textbf{tecweb.net} \quad Root: \code{/var/www/tw}};
    \node[f] (i1) at (-2.6,1.75) {index.html};
    \node[f, right=4pt of i1] (l1) {list.html};
    \node[f, right=4pt of l1] (e1) {example.html};
    \begin{scope}[on background layer]
      \node[host, fit=(h1t)(i1)(l1)(e1), inner sep=8pt] (h1) {};
    \end{scope}

    % --- host 2 ---
    \node[font=\Lato\scriptsize, anchor=west, text=primary!45!black] (h2t)
      at (-5.1,0.35) {Host: \textbf{bookofprogramming.com} \quad Root: \code{/var/www/book}};
    \node[f] (i2) at (-2.9,-0.35) {index.html};
    \node[d, right=4pt of i2] (p2) {pics/};
    \node[f, right=4pt of p2] (j2) {pics/fu-tzu.jpg};
    \begin{scope}[on background layer]
      \node[host, fit=(h2t)(i2)(p2)(j2), inner sep=8pt] (h2) {};
    \end{scope}

    \begin{scope}[on background layer]
      \node[rounded corners=5pt, draw=primary!70!black, line width=1pt,
            fill=primary!2!white, fit=(t)(h1)(h2), inner sep=10pt] {};
    \end{scope}
  \end{tikzpicture}
  \caption{Un solo server HTTP può ospitare più host, ciascuno con la propria
           document root.}
  \label{fig:http-server}
\end{figure}

\warning{Solo ciò che sta nella document root}{
  Il server serve esclusivamente i file contenuti nella document root. Tutto il
  resto del filesystem (il codice sorgente dell'applicazione, i file di
  configurazione con le password del database, le chiavi private) resta
  invisibile dall'esterno. È una regola semplice ma fondamentale: quando si
  pubblica un sito, \textbf{tutto ciò che finisce nella document root è
  pubblico}, compresi i file che pensavamo nessuno avrebbe indovinato.}

\section{Che cos'è HTML}

\dfn{HTML}{
  \term{HTML} (\textit{HyperText Markup Language}) è il linguaggio con cui si
  descrivono i documenti che i browser visualizzano. Una web app è costituita da
  uno o più documenti, detti anche \textit{pagine web}.}

Il concetto chiave è quello di \term{linguaggio di marcatura} (\textit{markup
language}): i documenti vengono \textbf{arricchiti} con un insieme di
annotazioni che ne controllano la struttura, la formattazione o la relazione fra
le loro parti.

Detto altrimenti: in un file HTML convivono due cose, il \textbf{contenuto} (il
testo che l'utente leggerà) e la \textbf{marcatura} (le annotazioni che dicono al
browser che cosa quel testo rappresenta: un titolo, un paragrafo, un
collegamento).

\subsection{Un po' di storia}

HTML ha attraversato diverse versioni dal 1993 a oggi. In questo corso ci
concentreremo sull'\term{HTML Living Standard}.

\begin{figure}[H]
  \centering
  \begin{tikzpicture}[font=\Lato\footnotesize,
      ver/.style={rounded corners=3pt, draw=primary!60!black, line width=0.8pt,
                  fill=primary!8!white, inner sep=4pt, align=center,
                  font=\Lato\scriptsize},
      cur/.style={ver, draw=accent!65!black, fill=accent!12!white,
                  font=\Lato\scriptsize\bfseries},
      tick/.style={draw=neutral!45!white, line width=0.6pt}]

    \draw[-{Stealth[length=7pt]}, line width=1pt, draw=neutral!45!black]
      (-0.7,0) -- (13.6,0);

    \foreach \x/\lbl in {0/{HTML 1.0\\1993}, 4.3/{HTML 3\\1997}, 8.6/{HTML 5\\2014}} {
      \draw[tick] (\x,0) -- (\x,0.45);
      \node[ver, anchor=south] at (\x,0.5) {\lbl};
      \fill[primary!70!black] (\x,0) circle (2.2pt);
    }
    \foreach \x/\lbl in {2.15/{HTML 2\\1995}, 6.45/{HTML 4\\1999}} {
      \draw[tick] (\x,0) -- (\x,-0.45);
      \node[ver, anchor=north] at (\x,-0.5) {\lbl};
      \fill[primary!70!black] (\x,0) circle (2.2pt);
    }

    \draw[tick] (11.6,0) -- (11.6,-0.45);
    \node[cur, anchor=north] at (11.6,-0.5) {HTML Living Standard\\dal 2019};
    \fill[accent!70!black] (11.6,0) circle (2.6pt);
  \end{tikzpicture}
  \caption{Le versioni di HTML.}
  \label{fig:html-timeline}
\end{figure}

\info{Che cosa significa \guillemotleft Living Standard\guillemotright}{
  Fino a HTML5 lo standard veniva pubblicato \textit{a versioni}: si congelava
  una specifica, le si dava un numero e si ricominciava. Dal 2019 il WHATWG ha
  adottato il modello di \textit{standard vivente}: esiste una sola specifica,
  sempre valida, aggiornata in continuazione. Non esiste (e non esisterà) un
  \guillemotleft HTML 6\guillemotright. Questo è il motivo per cui conviene consultare sempre la
  documentazione online e non un libro stampato.}

\section{Tag, elementi e attributi}

In HTML le annotazioni si chiamano \term{tag} e si scrivono racchiuse fra
parentesi angolari.

\begin{figure}[H]
  \centering
  \begin{tikzpicture}[
      part/.style={font=\FiraCode\large, inner sep=1pt, outer sep=0pt},
      lab/.style={font=\Lato\scriptsize, align=center, inner sep=2pt}]
    \node[part, text=codetag]                          (op) {<tagName>};
    \node[part, text=neutral!20!black, right=5pt of op] (ct) {contenuto del tag};
    \node[part, text=codetag, right=5pt of ct]         (cl) {</tagName>};

    \node[lab, text=codetag,         above=17pt of op] (lop) {Tag di apertura};
    \node[lab, text=neutral!20!black, above=17pt of ct] (lct) {Contenuto\\(opzionale)};
    \node[lab, text=codetag,         above=17pt of cl] (lcl) {Tag di chiusura};
    \draw[wtdash] (lop.south) -- (op.north);
    \draw[wtdash] (lct.south) -- (ct.north);
    \draw[wtdash] (lcl.south) -- (cl.north);

    \draw[decorate, decoration={brace, amplitude=5pt, mirror},
          draw=accent!70!black, line width=0.8pt]
      ([yshift=-5pt]op.south west) -- ([yshift=-5pt]cl.south east);
    \node[lab, text=accent!70!black, font=\Lato\small\bfseries]
      at ($(op.south west)!0.5!(cl.south east) + (0,-0.75)$) {Elemento HTML};
  \end{tikzpicture}
  \caption{Anatomia di un elemento HTML.}
  \label{fig:tag}
\end{figure}

Il tag di apertura può contenere anche \term{attributi}, coppie
\textbf{chiave-valore} (il valore è opzionale):

\begin{center}
  \code{<tagName attributo1="valore" attributo2>}
\end{center}

\subsection{Struttura di un documento}

Un documento HTML deve iniziare con la dichiarazione \code{<!DOCTYPE HTML>}. Non
è un tag: serve soltanto a comunicare al client quale tipo di documento
aspettarsi. Il tag \tg{html} rappresenta l'intero documento e contiene un
\tg{head} e un \tg{body}.

\begin{esempio}{Lo scheletro minimo di una pagina}
\begin{affianco}[0.5]
\begin{lstlisting}[style=html, name=hello.html]
<!DOCTYPE HTML>
<html lang="it">
  <head>
    <title>Ciao mondo!</title>
  </head>
  <body>
    <!-- questo e' un commento -->
    <p>Ciao mondo!</p>
  </body>
</html>
\end{lstlisting}
\risultato
\begin{browser}[Ciao mondo!]
  Ciao mondo!
\end{browser}
\wtnota{Il contenuto di \code{<title>} compare nella linguetta del browser,
        non nella pagina. Il commento non compare affatto.}
\end{affianco}

\begin{itemize}
  \item l'attributo \attr{lang} è fortemente consigliato per motivi di
        \textbf{accessibilità}: dice al browser e ai lettori di schermo in che
        lingua è scritto il contenuto, informazione fondamentale per la sintesi
        vocale e per la sillabazione;
  \item \tg{head} contiene i \textbf{metadati} del documento;
  \item \tg{body} contiene il \textbf{contenuto vero e proprio}, cioè quello che
        l'utente vedrà.
\end{itemize}
\end{esempio}

\section{L'elemento head}

L'elemento \tg{head} è un contenitore di \term{metadati}, cioè \guillemotleft dati sui
dati\guillemotright: informazioni sul documento corrente. Spesso non vengono mostrati
all'utente, ma sono utilissimi ai browser e ai motori di ricerca. È obbligatorio
che contenga un \tg{title}.

\begin{esempio}{Un head realistico}
\begin{lstlisting}[style=html, name=head.html]
<head>
  <title>Il Libro della Programmazione</title>
  <meta charset="UTF-8">
  <meta name="description" content="Frammenti dal Libro della Programmazione">
  <meta name="keywords" content="saggezza, programmazione">
  <meta name="author" content="L. L. L. Starace">
  <meta http-equiv="refresh" content="30">
</head>
\end{lstlisting}

\begin{itemize}
  \item \tg{title} è il titolo della pagina: compare nella linguetta del browser,
        nei segnalibri e come titolo nei risultati dei motori di ricerca;
  \item \attr{charset} dichiara la \textbf{codifica dei caratteri}, praticamente
        sempre \code{UTF-8};
  \item \attr{description} e \attr{keywords} sono metadati usati (con qualche
        cautela) dai motori di ricerca;
  \item \attr{http-equiv="refresh"} chiede al browser di ricaricare la pagina
        ogni 30 secondi.
\end{itemize}
\end{esempio}

\warning{Se manca il charset}{
  Se si dimentica \code{<meta charset="UTF-8">}, il browser deve
  \textit{indovinare} la codifica. Il risultato tipico è la comparsa di
  caratteri strani al posto delle lettere accentate: un problema che in italiano
  si nota immediatamente. La riga del charset va messa \textbf{il più in alto
  possibile} dentro \tg{head}.}

\subsection{La favicon}

Accanto al titolo, nella linguetta del browser, compare quasi sempre una piccola
icona: è la \term{favicon} (da \textit{favorite icon}), che il browser mostra
anche nei segnalibri, nella cronologia e nell'elenco delle schede aperte. Si
dichiara nell'\tg{head} con un elemento \tg{link}:

\begin{affianco}[0.5]
\begin{lstlisting}[style=html, name=index.html]
<head>
  <meta charset="UTF-8">
  <title>Il Libro della Programmazione</title>
  <link rel="icon" href="/favicon.ico">
</head>
\end{lstlisting}
\risultato
\begin{center}
\begin{tikzpicture}[font=\Lato\scriptsize]
  % barra delle schede
  \fill[neutral!12!white, rounded corners=3pt] (0,0) rectangle (7.6,1.55);
  \fill[white, rounded corners=3pt] (0.15,0.78) rectangle (5.3,1.42);
  \node[anchor=west] (ico) at (0.3,1.1)
    {\tikz{\fill[primary!80!black, rounded corners=1pt] (0,0) rectangle (0.28,0.28);
           \node[text=white, font=\Lato\bfseries\tiny, inner sep=0pt] at (0.14,0.14) {F};}};
  \node[anchor=west, text=neutral!15!black] (tit) at (0.75,1.1)
    {Il Libro della Programmazione};
  \node[text=neutral!45!black] at (5.05,1.1) {$\times$};
  % barra dell'indirizzo
  \fill[white] (0,0) rectangle (7.6,0.72);
  \draw[neutral!35!white, rounded corners=5pt] (0.9,0.14) rectangle (7.4,0.58);
  \node[text=neutral!45!black] at (0.28,0.36) {$\leftarrow$};
  \node[anchor=west, text=neutral!40!black, font=\FiraCodeNL\tiny] at (1.05,0.36)
    {127.0.0.1:3000/index.html};
  \draw[neutral!35!white] (0,0) rectangle (7.6,1.55);
  % etichette
  \draw[wtarrowacc] (0.46,2.2) -- (0.46,1.35);
  \node[wtlab, anchor=south] at (0.46,2.2) {favicon\\\code{<link rel="icon">}};
  \draw[wtarrow] (3.2,2.2) -- (3.2,1.35);
  \node[wtlab, anchor=south] at (3.2,2.2) {titolo\\\code{<title>}};
\end{tikzpicture}
\end{center}
\end{affianco}

\begin{itemize}
  \item \attr{rel="icon"} dice che la risorsa collegata è l'\textbf{icona} del
        documento (lo stesso attributo \attr{rel} che, con valore
        \code{stylesheet}, vedremo servire a collegare i fogli di stile);
  \item \attr{href} indica la URL dell'immagine. Il formato storico è
        \code{.ico}, ma oggi si usano senza problemi anche \code{.png} e
        \code{.svg}, eventualmente indicando il tipo con
        \code{type="image/png"}.
\end{itemize}

La cosa più interessante è ciò che accade dietro le quinte: la favicon
\textbf{non è contenuta} nel documento HTML, quindi il browser deve effettuare
una \textbf{richiesta HTTP aggiuntiva} per scaricarla, esattamente come vedremo
per le immagini. C'è di più: anche se la pagina \textit{non} dichiara alcuna
icona, la maggior parte dei browser prova comunque a chiedere al server il file
\code{/favicon.ico} nella radice del sito. Se il file non esiste, il server
risponde \code{404 Not Found}: la pagina funziona lo stesso, ma quella richiesta
fallita compare nei log del server e nella scheda \textit{Rete} dei DevTools,
come vedremo alla fine del capitolo (figura~\ref{fig:devtools-rete}).

\info{Perché la favicon «non si aggiorna»}{
  I browser conservano le favicon nella propria \textbf{cache} in modo molto
  tenace: sostituito il file sul server, capita di continuare a vedere per un
  po' l'icona vecchia. Di solito basta una ricarica forzata della pagina
  (\code{Ctrl+Shift+R}) o aprire direttamente la URL dell'icona. Del
  funzionamento della cache parliamo alla fine del capitolo.}

\section{Elementi fondamentali}

\subsection{Intestazioni}

Gli elementi da \code{<h1>} a \code{<h6>} rappresentano \term{titoli} e
sottotitoli, dal più importante al meno importante.

\begin{esempio}{I sei livelli di intestazione}
\begin{affianco}[0.46]
\begin{lstlisting}[style=html, name=headings.html]
<body>
  <h1>Titolo di primo livello</h1>
  <h2>Titolo di secondo livello</h2>
  <h3>Titolo di terzo livello</h3>
  <h4>Titolo di quarto livello</h4>
  <h5>Titolo di quinto livello</h5>
  <h6>Titolo di sesto livello</h6>
</body>
\end{lstlisting}
\risultato
\begin{browser}[headings.html]
  {\Lora\bfseries\Huge Titolo di primo livello}\\[5pt]
  {\Lora\bfseries\LARGE Titolo di secondo livello}\\[5pt]
  {\Lora\bfseries\large Titolo di terzo livello}\\[5pt]
  {\Lora\bfseries\normalsize Titolo di quarto livello}\\[5pt]
  {\Lora\bfseries\small Titolo di quinto livello}\\[5pt]
  {\Lora\bfseries\footnotesize Titolo di sesto livello}
\end{browser}
\end{affianco}
\end{esempio}

\warning{Le intestazioni non servono a fare il testo grande}{
  Un errore molto comune fra i principianti è scegliere \code{<h3>} invece di
  \code{<h1>} \guillemotleft perché la dimensione mi piace di più\guillemotright. Le intestazioni
  esprimono la \textbf{gerarchia} del documento, non il suo aspetto: i motori di
  ricerca le usano per capire di che cosa parla la pagina, e i lettori di schermo
  per permettere agli utenti non vedenti di navigarla saltando di sezione in
  sezione. Per cambiare la dimensione del testo esiste il CSS.}

\subsection{Paragrafi}

L'elemento \tg{p} rappresenta un \term{paragrafo}. I paragrafi iniziano
tipicamente su una nuova riga.

\begin{esempio}{Due paragrafi}
\begin{affianco}[0.46]
\begin{lstlisting}[style=html, name=paragrafi.html]
<h1>Il Libro della Programmazione</h1>
<p>
  Questa pagina contiene alcuni frammenti
  del grande Libro della Programmazione,
  un antico manoscritto perduto nel tempo.
</p>
<p>
  Si dice che il libro sia stato scritto da
  uno degli studenti del leggendario maestro
  Fu-Tzu.
</p>
\end{lstlisting}
\risultato
\begin{browser}[paragrafi.html]
  {\Lora\bfseries\LARGE Il Libro della Programmazione}\\[6pt]
  Questa pagina contiene alcuni frammenti del grande Libro della Programmazione,
  un antico manoscritto perduto nel tempo.\\[6pt]
  Si dice che il libro sia stato scritto da uno degli studenti del leggendario
  maestro Fu-Tzu.
\end{browser}
\end{affianco}

Si noti una cosa importante: nel sorgente il testo va a capo più volte e
contiene molti spazi di indentazione, ma nel risultato tutto ciò è
\textbf{scomparso}. HTML \term{collassa gli spazi bianchi}: qualunque sequenza di
spazi, tabulazioni e a-capo consecutivi viene trattata come un \textbf{unico
spazio}. Se serve andare a capo davvero, occorre dirlo esplicitamente con
\tg{br}.
\end{esempio}

\subsection{Commenti}

I \term{commenti} sono delimitati da \code{<!-{}-} e \code{-{}->} e vengono
ignorati dai browser. Sono utili per aggiungere note o per nascondere
temporaneamente del contenuto.

\begin{esempio}{Commenti singoli e multi-riga}
\begin{affianco}[0.54]
\begin{lstlisting}[style=html, name=commenti.html]
<h1>I due aspetti</h1>
<!-- TODO: aggiungere altra saggezza -->
<p>
  Sotto la superficie della macchina,
  il programma si muove. Senza sforzo,
  si espande e si contrae.
</p>
<!--
<p>
  Le forme sul monitor non sono che
  increspature sull'acqua. L'essenza
  resta invisibile, al di sotto.
</p> -->
\end{lstlisting}
\risultato
\begin{browser}[commenti.html]
  {\Lora\bfseries\Huge I due aspetti}\\[6pt]
  Sotto la superficie della macchina, il programma si muove. Senza sforzo, si
  espande e si contrae.
\end{browser}
\wtnota{Il secondo paragrafo è dentro un commento: non viene mostrato.}
\end{affianco}
\end{esempio}

\error{I commenti non sono un nascondiglio}{
  Il contenuto commentato \textbf{non} viene mostrato, ma viene comunque
  \textbf{inviato al browser}: basta guardare il sorgente della pagina per
  leggerlo. Non commentare mai credenziali, chiavi API, note interne o indirizzi
  di ambienti di test: chiunque può leggerli con due clic.}

\subsection{Semantica del testo}

Alcuni tag specificano la \term{semantica} di porzioni di testo.

\begin{center}
\renewcommand{\arraystretch}{1.35}
\begin{tabular}{@{}l@{\hspace{1.1cm}}>{\raggedright\arraybackslash}p{9.4cm}@{}}
  \textbf{Tag} & \textbf{Significato} \\
  \code{<em>}     & Enfatizza il contenuto. \\
  \code{<strong>} & Indica \textit{forte importanza}. \\
  \code{<br/>}    & Interruzione di riga. È un \term{tag vuoto}, quindi può
                    essere auto-chiuso. \\
  \code{<abbr>}   & Definisce acronimi e abbreviazioni; il testo esteso va
                    nell'attributo \attr{title}. \\
  \code{<del>}    & Contenuto che è stato \textit{rimosso} dal documento. \\
  \code{<ins>}    & Contenuto che è stato \textit{inserito} nel documento. \\
\end{tabular}
\end{center}

\begin{esempio}{Semantica testuale in azione}
\begin{affianco}[0.55]
\begin{lstlisting}[style=html, name=semantica.html]
<p>
  I tag
  <abbr title="HyperText Markup Language">HTML</abbr>
  possono essere usati con semantica a livello
  di testo.<br/> Alcuni tag conferiscono
  <strong>importanza</strong> al loro contenuto,
  altri lo <em>enfatizzano</em>. Alcuni servono
  a evidenziare contenuto <del>rimosso</del>
  <ins>aggiunto</ins>.
</p>
\end{lstlisting}
\risultato
\begin{browser}[semantica.html]
  I tag \dotuline{HTML} possono essere usati con semantica a livello di testo.\\
  Alcuni tag conferiscono \textbf{importanza} al loro contenuto, altri lo
  \textit{enfatizzano}. Alcuni servono a evidenziare contenuto \sout{rimosso}
  \uline{aggiunto}.
\end{browser}
\end{affianco}

Il testo esteso dell'abbreviazione (\textit{HyperText Markup Language}) compare
come suggerimento quando il puntatore si ferma sopra la sigla, che il browser
sottolinea con una linea punteggiata.
\end{esempio}

\info{Semantica, non aspetto}{
  Esistono anche \code{<i>} (corsivo) e \code{<b>} (grassetto), che descrivono
  soltanto l'\textit{aspetto}. Si preferiscono \code{<em>} e \code{<strong>}
  perché descrivono il \textit{significato}: un lettore di schermo cambia
  intonazione su \code{<strong>}, mentre su \code{<b>} non ha nulla da dire. La
  regola generale in HTML è sempre la stessa: \textbf{descrivi che cosa è, non
  come deve apparire}. Come deve apparire lo dirà il CSS.}

\section{Collegamenti ipertestuali}

I collegamenti si definiscono con il tag \term{anchor} \tg{a}. L'attributo
\attr{href} indica la URL di destinazione.

\begin{esempio}{Un collegamento}
\begin{affianco}[0.5]
\begin{lstlisting}[style=html, name=link.html]
<p>
  <a href="/fu-tzu.html">Fu-Tzu</a> una volta disse:
  <em>
    Quando il saggio indica la luna,
    lo sciocco fissa il dito e si chiede
    da quale libreria provenga.
  </em>
</p>
\end{lstlisting}
\risultato
\begin{browser}[link.html]
  \wtlink{Fu-Tzu} una volta disse: \textit{Quando il saggio
  indica la luna, lo sciocco fissa il dito e si chiede da quale libreria
  provenga.}
\end{browser}
\end{affianco}
\end{esempio}

\subsection{URL relative e assolute}

Le URL indicate in \attr{href} possono essere \term{assolute} o \term{relative}.

\begin{itemize}
  \item Le \textbf{URL assolute} includono schema e nome di host, e contengono
        tutte le informazioni necessarie a raggiungere la risorsa: \\
        \code{<a href="http://tecweb.com/course/intro.html">}
  \item Le \textbf{URL relative} indicano solo un percorso; schema e host vengono
        dedotti dal contesto corrente: \\
        \code{<a href="page.html">} \quad oppure \quad
        \code{<a href="/course/react.html">}
\end{itemize}

La regola di risoluzione dipende dal fatto che la URL relativa cominci o meno
con una barra.

\begin{esempio}{URL relativa che inizia con la barra}
Quando una URL relativa inizia con \code{/}, l'\textbf{intero percorso} viene
sostituito.

\begin{itemize}
  \item Contesto corrente:
        \code{http://bookofprogramming.com/fu-tzu/fu-tzu.html}
  \item Un'ancora \code{<a href="/index.html">} punta a:
        \code{http://bookofprogramming.com/index.html}
\end{itemize}
\end{esempio}

\begin{esempio}{URL relativa che non inizia con la barra}
Quando una URL relativa \textbf{non} inizia con \code{/}, viene sostituito solo
l'\textbf{ultimo segmento} del percorso, cioè quello dopo l'ultima barra (qui,
il nome del file).

\begin{itemize}
  \item Contesto corrente:
        \code{http://bookofprogramming.com/fu-tzu/fu-tzu.html}
  \item Un'ancora \code{<a href="pic.jpg">} punta a:
        \code{http://bookofprogramming.com/fu-tzu/pic.jpg}
\end{itemize}
\end{esempio}

\subsection{I segmenti punto}

Le URL relative possono contenere anche i cosiddetti \term{segmenti punto}:

\begin{itemize}
  \item il punto singolo (\code{.}) rappresenta la \textbf{directory corrente};
  \item il punto doppio (\code{..}) rappresenta la \textbf{directory padre}.
\end{itemize}

Supponendo che il percorso corrente sia \code{/a/b/c/hello.html}:

\begin{center}
\renewcommand{\arraystretch}{1.35}
\begin{tabular}{@{}l@{\hspace{1.5cm}}l@{}}
  \textbf{href} & \textbf{Percorso risultante} \\
  \code{./index.html}    & \code{/a/b/c/index.html} \\
  \code{../foo.html}     & \code{/a/b/foo.html} \\
  \code{../../pic.jpg}   & \code{/a/pic.jpg} \\
  \code{./../../pic.jpg} & \code{/a/pic.jpg} \quad (identico al precedente) \\
\end{tabular}
\end{center}

\subsection{Meglio relative o assolute?}

\begin{itemize}
  \item Le \textbf{URL relative} vanno preferite quando si collegano risorse
        \textbf{interne} alla stessa applicazione web. In questo modo, se cambia
        il nome di host (per esempio passando dall'ambiente di sviluppo a quello
        di produzione), non c'è alcun bisogno di modificare le pagine HTML.
  \item Per collegare pagine o risorse \textbf{esterne} non c'è scelta: si devono
        usare URL assolute.
\end{itemize}

\subsection{L'attributo target}

L'attributo \attr{target} specifica \textit{dove} aprire la risorsa collegata:

\begin{itemize}
  \item \code{<a target="\_self" href="pic.jpg">} apre la risorsa nella stessa
        finestra o scheda del browser (è il comportamento predefinito, quando
        \attr{target} non viene specificato);
  \item \code{<a target="\_blank" href="pic.jpg">} apre la risorsa in una
        \textbf{nuova} finestra o scheda.
\end{itemize}

\section{URL che puntano a una directory}

Che cosa succede quando una URL punta a una \textbf{directory} e non a un file?

I server web rispondono tipicamente con il contenuto del file \code{index.html},
se esiste nella directory richiesta. Questo comportamento è configurabile:

\begin{itemize}
  \item altri nomi predefiniti comunemente usati sono \code{home.html} e
        \code{default.html};
  \item se non esiste alcun file predefinito, i server HTTP possono essere
        configurati per generare automaticamente un \term{indice} della
        directory, cioè un elenco navigabile dei file che contiene.
\end{itemize}

\warning{L'indice automatico è comodo e pericoloso}{
  L'elenco automatico dei file è utilissimo in sviluppo, ma su un server di
  produzione espone la struttura interna del sito: backup dimenticati, archivi
  \code{.zip}, vecchie versioni delle pagine. È buona norma disattivarlo e
  mettere sempre un \code{index.html}, anche vuoto.}

\section{Tabelle}

Una \term{tabella} \tg{table} contiene un insieme di righe \tg{tr}
(\textit{table row}). Ogni riga può contenere una o più celle:

\begin{itemize}
  \item \tg{td} (\textit{table data}), le celle di dati;
  \item \tg{th} (\textit{table header}), le celle di intestazione.
\end{itemize}

Una tabella può contenere anche un elemento \tg{caption} che la descrive.

\begin{esempio}{Una tabella di esami}
\begin{affianco}[0.5]
\begin{lstlisting}[style=html, name=tabella.html]
<table>
  <caption>Esami</caption>
  <tr>
    <th>Esame</th><th>Voto</th>
  </tr>
  <tr>
    <td>A</td><td>30</td>
  </tr>
  <tr>
    <td>B</td><td>28</td>
  </tr>
</table>
\end{lstlisting}
\risultato
\begin{browser}[tabella.html]
  \begin{center}
    {\Lora Esami}\\[3pt]
    {\Lora
    \begin{tabular}{|l|l|}
      \hline
      \textbf{Esame} & \textbf{Voto} \\ \hline
      A & 30 \\ \hline
      B & 28 \\ \hline
    \end{tabular}}
  \end{center}
\end{browser}
\end{affianco}

Senza CSS il browser disegna la tabella in modo molto spartano: nessun bordo
(nell'esempio i bordi sono aggiunti per chiarezza), intestazioni in grassetto e
centrate, celle larghe quanto basta.
\end{esempio}

\error{Le tabelle non servono a impaginare}{
  Negli anni Novanta e Duemila era prassi comune costruire l'\textit{intero
  layout} di un sito con tabelle annidate. Oggi è un grave errore: rende la
  pagina incomprensibile ai lettori di schermo, impossibile da rendere
  \textit{responsive} e difficilissima da manutenere. Le tabelle servono a
  presentare \textbf{dati tabellari}. Per l'impaginazione esistono i CSS, che
  vedremo nel prossimo capitolo.}

\section{Liste}

HTML definisce tre tipi di lista:

\begin{itemize}
  \item liste \term{ordinate} \tg{ol}, usate per le enumerazioni;
  \item liste \term{non ordinate} \tg{ul}, usate per gli elenchi puntati;
  \item liste di \term{descrizione} \tg{dl}, costituite da termini e relative
        descrizioni, spesso usate per i glossari.
\end{itemize}

\subsection{Liste ordinate e non ordinate}

Le liste ordinate e non ordinate contengono una sequenza di elementi di lista
\tg{li} (\textit{list item}).

\begin{esempio}{Le due liste a confronto}
\begin{affianco}[0.5]
\begin{lstlisting}[style=html, name=liste.html]
<ol>
  <li>Anna</li>
  <li>Bruno</li>
  <li>Carlo</li>
</ol>

<ul>
  <li>TecWeb</li>
  <li><a href="/">Fu-Tzu</a></li>
</ul>
\end{lstlisting}
\risultato
\begin{browser}[liste.html]
  \begin{enumerate}[leftmargin=1.8em, itemsep=0pt, topsep=2pt, parsep=0pt]
    \item Anna
    \item Bruno
    \item Carlo
  \end{enumerate}
  \vspace{4pt}
  \begin{itemize}[leftmargin=1.8em, itemsep=0pt, topsep=2pt, parsep=0pt]
    \item TecWeb
    \item \wtlink{Fu-Tzu}
  \end{itemize}
\end{browser}
\end{affianco}
\end{esempio}

\subsection{Liste di descrizione}

Le liste di descrizione contengono una sequenza di termini \tg{dt}
(\textit{description term}) e di descrizioni \tg{dd} (\textit{description
details}) relative ai termini che le precedono.

\begin{esempio}{Un piccolo glossario}
\begin{affianco}[0.5]
\begin{lstlisting}[style=html, name=glossario.html]
<dl>
  <dt>Dati</dt>
  <dd>
    Sono fatti di bit,
    eppure assumono forme complesse.
  </dd>
  <dt>Controllo</dt>
  <dd>
    Consiste di istruzioni semplici,
    eppure svolge compiti difficili.
  </dd>
</dl>
\end{lstlisting}
\risultato
\begin{browser}[glossario.html]
  \textbf{Dati}\\
  \hspace*{2em}Sono fatti di bit, eppure assumono forme complesse.\\[3pt]
  \textbf{Controllo}\\
  \hspace*{2em}Consiste di istruzioni semplici, eppure svolge compiti difficili.
\end{browser}
\end{affianco}
\end{esempio}

\section{Riferimenti a caratteri (entità)}

Alcuni caratteri sono \term{riservati} in HTML. Che cosa succede se vogliamo
scrivere in un documento \code{<p>3<x e y>6</p>}?

Il browser rischia di confondere i simboli con i tag: incontrando \code{<x} si
aspetta l'inizio di un tag chiamato \code{x}. Per visualizzare i caratteri
riservati bisogna usare i \term{riferimenti a caratteri} (\textit{character
references}, chiamati spesso anche \term{entità}), che si scrivono in due modi:

\begin{center}
  \code{\&nome\_entita;} \qquad oppure \qquad \code{\&\#numero\_entita;}
\end{center}

\begin{center}
\renewcommand{\arraystretch}{1.35}
\begin{tabular}{@{}c@{\hspace{1.2cm}}l@{\hspace{1.2cm}}l@{}}
  \textbf{Risultato} & \textbf{Descrizione} & \textbf{Riferimento} \\
             & Spazio non divisibile           & \code{\&nbsp;}  \\
  \code{<}   & Minore di                       & \code{\&lt;}    \\
  \code{>}   & Maggiore di                     & \code{\&gt;}    \\
  \code{\&}  & E commerciale                   & \code{\&amp;}   \\
  \code{"}   & Virgolette doppie               & \code{\&quot;}  \\
  \code{'}   & Virgolette singole (apostrofo)  & \code{\&apos;}  \\
  \copyright & Copyright                       & \code{\&copy;}  \\
\end{tabular}
\end{center}

L'esempio visto sopra andrebbe quindi scritto così:

\begin{affianco}[0.5]
\begin{lstlisting}[style=html, name=riferimenti.html]
<p>3&lt;x e y&gt;6</p>
<p>3&#60;x e y&#62;6</p>
\end{lstlisting}
\risultato
\begin{browser}[riferimenti.html]
  3<x e y>6\\[4pt]
  3<x e y>6
\end{browser}
\end{affianco}

Le due righe sono equivalenti: \code{\&lt;} è il riferimento per
\textbf{nome} e \code{\&\#60;} quello per \textbf{numero} (60 è il codice del
carattere \code{<}). La tabella completa dei riferimenti si trova nella
specifica HTML.

\info{Lo spazio non divisibile}{
  \code{\&nbsp;} (\textit{non-breaking space}) è particolarmente utile: essendo
  un carattere \guillemotleft vero\guillemotright\ non viene collassato con gli spazi adiacenti e
  impedisce al browser di andare a capo in quel punto. Serve, per esempio, a
  tenere unite espressioni come \guillemotleft 10 kg\guillemotright\ o \guillemotleft Fig. 3\guillemotright.}

\section{Immagini}

L'elemento \tg{img} serve a incorporare un'immagine in un documento HTML.

\begin{itemize}
  \item È un \term{elemento vuoto}: non deve contenere altri elementi;
  \item l'attributo \attr{src} specifica la URL dell'immagine da includere;
  \item l'attributo \attr{alt} specifica un testo alternativo che la descrive;
  \item gli attributi \attr{width} e \attr{height} permettono di specificarne le
        dimensioni in pixel all'interno della pagina.
\end{itemize}

\begin{esempio}{Inserire un'immagine}
\begin{affianco}[0.46]
\begin{lstlisting}[style=html, name=image.html]
<h1>Fu-Tzu</h1>
<p>
  Un ritratto del maestro,
  generato da DALL-E, e' disponibile
  qui di seguito:
</p>
<img src="pic.jpg"
     alt="ritratto di Fu-Tzu"
     width="200"/>
\end{lstlisting}
\risultato
\begin{browser}[image.html]
  {\Lora\bfseries\Huge Fu-Tzu}\\[6pt]
  Un ritratto del maestro, generato da DALL-E, è disponibile qui di
  seguito:\\[6pt]
  \includegraphics[width=3.4cm]{sources/fu-tzu.jpg}
\end{browser}
\end{affianco}
\end{esempio}

\warning{L'attributo alt non è facoltativo}{
  Il testo alternativo viene letto dai lettori di schermo agli utenti non
  vedenti, viene mostrato se l'immagine non si carica ed è usato dai motori di
  ricerca per capire che cosa raffigura. Ometterlo significa rendere una parte
  del contenuto invisibile a una parte degli utenti. Fa eccezione il caso delle
  immagini puramente decorative, per le quali si usa \code{alt=""} (vuoto), che
  dice esplicitamente ai lettori di schermo di ignorarle.}

\subsection{Dietro le quinte}

Con le immagini succede qualcosa di nuovo, e piuttosto interessante.

Fino a questo momento i documenti HTML erano \textbf{autocontenuti}: tutti i dati
del documento stavano \textit{dentro} il documento. C'erano dei collegamenti, ma
eravamo liberi di non seguirli.

L'ultima pagina d'esempio, invece, include un contenuto esterno (l'immagine)
fornendone soltanto la URL. L'immagine \textbf{non è} inclusa nel documento HTML:
per visualizzare la pagina il browser deve procurarsi risorse aggiuntive.

\begin{figure}[H]
  \centering
  \begin{tikzpicture}[font=\Lato\small,
      hdr/.style={wtnode, minimum width=3.0cm, minimum height=0.9cm},
      life/.style={draw=neutral!45!white, line width=0.7pt, dashed},
      msg/.style={-{Stealth[length=6pt,width=5pt]}, line width=0.9pt,
                  draw=primary!70!black},
      rsp/.style={-{Stealth[length=6pt,width=5pt]}, line width=0.9pt,
                  draw=accent!70!black},
      ml/.style={font=\Lato\scriptsize, midway, above, inner sep=2.5pt,
                 fill=white, text=neutral!15!black},
      act/.style={font=\Lato\scriptsize\itshape, text=neutral!30!black,
                  align=left, anchor=west},
      num/.style={circle, fill=accent!75!black, text=white, inner sep=1.5pt,
                  font=\Lato\bfseries\tiny}]

    \def\xb{0} \def\xs{9.6}
    \node[hdr] (B) at (\xb,0) {\textbf{Browser}};
    \node[hdr] (S) at (\xs,0) {\textbf{Server HTTP}};
    \draw[life] (B.south) -- (\xb,-5.5);
    \draw[life] (S.south) -- (\xs,-5.5);

    \draw[msg] (\xb,-0.9) -- (\xs,-0.9) node[ml]{\code{GET /02-HTML/image.html}};
    \node[num] at (-0.5,-0.9) {1};

    \draw[rsp] (\xs,-1.8) -- (\xb,-1.8) node[ml]{\code{200 OK} + documento HTML};
    \node[num] at (-0.5,-1.8) {2};

    \node[act] at (0.25,-2.7)
      {il browser analizza il documento dall'alto verso il basso\\
       e incontra \code{<img src="pic.jpg">}};
    \node[num] at (-0.5,-2.7) {3};

    \draw[msg] (\xb,-3.7) -- (\xs,-3.7) node[ml]{\code{GET /02-HTML/pic.jpg}};
    \node[num] at (-0.5,-3.7) {4};

    \draw[rsp] (\xs,-4.6) -- (\xb,-4.6) node[ml]{\code{200 OK} + dati dell'immagine};
    \node[num] at (-0.5,-4.6) {5};

    \node[act] at (0.25,-5.3) {la pagina viene mostrata completa};
    \node[num] at (-0.5,-5.3) {6};
  \end{tikzpicture}
  \caption{Una singola pagina con un'immagine richiede due richieste HTTP.}
  \label{fig:img-fetch}
\end{figure}

\info{Una pagina, tante richieste}{
  Questo meccanismo si applica a \textbf{tutte} le risorse esterne: immagini,
  fogli di stile CSS, script JavaScript, font, video. Una pagina moderna può
  facilmente generare cinquanta o cento richieste HTTP. È per questo che il
  numero e la dimensione delle risorse esterne incidono così tanto sulla
  velocità percepita di un sito, ed è per questo che la scheda \textit{Network}
  dei DevTools è uno strumento diagnostico tanto importante.}

\section{Attributi globali}

Abbiamo incontrato diversi attributi (\attr{href}, \attr{target}, \attr{src},
\attr{alt}). Questi attributi hanno senso solo per alcuni elementi:
\code{<strong src="pic.jpg">Ciao!</strong>} non significherebbe nulla.

Alcuni attributi sono invece \term{globali}, cioè utilizzabili con
\textbf{qualunque} elemento HTML.

\begin{center}
\renewcommand{\arraystretch}{1.35}
\begin{tabular}{@{}l@{\hspace{1.1cm}}>{\raggedright\arraybackslash}p{9.4cm}@{}}
  \textbf{Attributo} & \textbf{Descrizione} \\
  \attr{id}    & Specifica un identificatore \textbf{univoco} nel documento. \\
  \attr{lang}  & Specifica la lingua del contenuto dell'elemento. \\
  \attr{style} e \attr{class} & Usati per la presentazione grafica (li vedremo
                                nel prossimo capitolo). \\
\end{tabular}
\end{center}

\subsection{Ancore nelle URL}

L'attributo \attr{id} può essere usato anche nelle \term{ancore} delle URL.

\begin{figure}[H]
  \centering
  \begin{tikzpicture}[
      part/.style={font=\FiraCode\large, inner sep=1pt, outer sep=0pt},
      lab/.style={font=\Lato\scriptsize, align=center, inner sep=2pt}]
    \node[part, text=codekey]                       (sch) {https://};
    \node[part, text=codetag,    right=0pt of sch]  (dom) {app.com};
    \node[part, text=codestring, right=0pt of dom]  (pth) {/foo/bar.html};
    \node[part, text=codeattr,   right=0pt of pth]  (anc) {\#some-id};

    \node[lab, text=codekey,  above=17pt of sch] (l1) {Schema};
    \node[lab, text=codeattr, above=17pt of anc] (l4) {Ancora};
    \draw[wtdash] (l1.south) -- (sch.north);
    \draw[wtdash] (l4.south) -- (anc.north);

    \node[lab, text=codetag,    below=17pt of dom] (l2) {Nome di dominio};
    \node[lab, text=codestring, below=17pt of pth] (l3) {Percorso};
    \draw[wtdash] (l2.north) -- (dom.south);
    \draw[wtdash] (l3.north) -- (pth.south);
  \end{tikzpicture}
  \caption{Un'ancora punta a un punto preciso \textit{dentro} la risorsa.}
  \label{fig:url-anchor}
\end{figure}

Un'ancora rappresenta una sorta di \guillemotleft segnalibro\guillemotright\ interno alla risorsa: dice
al browser di mostrare direttamente il contenuto che si trova in quel punto.
Nell'esempio, \code{\#some-id} punta all'elemento il cui attributo \attr{id} vale
\code{some-id}. Un caso reale è
\code{https://luistar.github.io/publications/\#conference}.

\section{Form}

Gli elementi \tg{form} servono a raccogliere \term{input dell'utente}. L'input
viene tipicamente inviato a un server per l'elaborazione.

I form contengono elementi come \tg{input}, \tg{label}, \tg{textarea} e
\tg{select}.

\begin{esempio}{Un form minimale}
\begin{affianco}[0.5]
\begin{lstlisting}[style=html, name=form.html]
<form>
  Messaggio:
  <input type="text"><br><br>
  <input type="submit">
</form>
\end{lstlisting}
\risultato
\begin{browser}[form.html]
  Messaggio: \wtcampo[3.5cm]{}\\[8pt]
  \wtbottone{Invia}
\end{browser}
\end{affianco}
Il pulsante \code{submit} non ha un'etichetta nel codice: il testo
«Invia» lo sceglie il browser, nella lingua dell'utente.
\end{esempio}

\subsection{Tipi di input}

Tipi diversi di elemento \tg{input} sono disponibili tramite l'attributo
\attr{type}.

\begin{center}
\renewcommand{\arraystretch}{1.35}
\begin{tabular}{@{}l@{\hspace{1.1cm}}>{\raggedright\arraybackslash}p{8.6cm}@{}}
  \textbf{Tipo} & \textbf{Descrizione} \\
  \code{text}     & Campo di testo su una sola riga. \\
  \code{password} & Campo per password (l'input è nascosto con asterischi). \\
  \code{number}   & Campo per numeri. \\
  \code{radio}    & Pulsante radio, per scegliere \textbf{una} fra molte opzioni. \\
  \code{checkbox} & Casella di spunta, per sceglierne \textbf{zero o più}. \\
  \code{button}   & Pulsante cliccabile. \\
\end{tabular}
\end{center}

\subsection{Etichette}

L'elemento \tg{label} serve a \term{etichettare} ciascun input. Usare le
etichette è una buona pratica di usabilità e accessibilità. L'attributo
\attr{for} della label deve essere uguale all'\attr{id} dell'input
corrispondente.

\begin{esempio}{Un form di login}
\begin{affianco}[0.55]
\begin{lstlisting}[style=html, name=login.html]
<form>
  <label for="id_usr">Nome utente: </label>
  <input type="text" name="usr" id="id_usr">
  <br/>
  <label for="id_pwd">Password: </label>
  <input type="password" name="pwd" id="id_pwd">
  <br/>
  <input type="submit">
</form>
\end{lstlisting}

\risultato
\begin{browser}[login.html]
  Nome utente: \wtcampo[3cm]{fu-tzu}\\[4pt]
  Password: \wtcampo[3cm]{\textbullet\textbullet\textbullet\textbullet\textbullet\textbullet}\\[8pt]
  \wtbottone{Invia}
\end{browser}
\end{affianco}

Oltre al vantaggio per i lettori di schermo, le etichette collegate rendono
cliccabile il testo stesso: fare clic su \guillemotleft Nome utente\guillemotright\ porta il cursore
dentro il campo. È un dettaglio che migliora molto l'usabilità, soprattutto su
schermi piccoli.
\end{esempio}

\subsection{Invio del form}

I form possono essere \term{inviati} (\textit{submit}) per trasmettere i dati
raccolti a un gestore (\textit{form-handler}). Al momento dell'invio viene
tipicamente effettuata una nuova richiesta HTTP.

\begin{itemize}
  \item La URL del gestore a cui la richiesta viene inviata è specificata
        dall'attributo \attr{action} del form (per impostazione predefinita, la
        stessa pagina su cui si trova il form).
  \item È possibile specificare anche il metodo HTTP da usare, tramite
        l'attributo \attr{method} (per impostazione predefinita, \code{GET}).
\end{itemize}

\begin{lstlisting}[style=html]
<form action="/handler.html" method="GET">
  <!-- qui gli elementi del form -->
</form>
\end{lstlisting}

Al momento dell'invio, l'input raccolto viene rappresentato come una serie di
coppie \textbf{nome/valore} nella forma:

\begin{center}
  \code{nome1=valore1\&...\&nomeN=valoreN}
\end{center}

Ogni \textit{nome} è il valore dell'attributo \attr{name} di un elemento di
input, e il \textit{valore} corrispondente è il suo contenuto al momento
dell'invio.

\warning{name e id non sono la stessa cosa}{
  L'attributo \attr{id} identifica l'elemento \textbf{dentro il documento} e
  serve al CSS, a JavaScript e alle etichette. L'attributo \attr{name} è il nome
  con cui il valore viene \textbf{inviato al server}. Un input privo di
  \attr{name} viene semplicemente \textbf{ignorato} al momento dell'invio: è uno
  degli errori più frequenti e più difficili da individuare, perché la pagina
  sembra funzionare perfettamente.}

\subsection{Invio con GET}

Se il metodo è \code{GET}, gli input vengono \textbf{accodati alla URL} del
gestore.

\begin{esempio}{Invio con il metodo GET}
\begin{affianco}[0.53]
\begin{lstlisting}[style=html, name=form-get.html]
<form action="/handler.html" method="GET">
  Messaggio: <input type="text" name="msg"><br>
  Numero: <input type="number" name="num"><br>
  <input type="submit">
</form>
\end{lstlisting}
\risultato
\begin{browser}[form-get.html]
  Messaggio: \wtcampo[2.6cm]{Hello!}\\[3pt]
  Numero: \wtcampo[2.6cm]{42}\\[5pt]
  \wtbottone{Invia}
\end{browser}
\wtnota{Premendo \textit{Invia} $\downarrow$}
\begin{lstlisting}[style=http, name=Richiesta generata]
GET /handler.html?msg=Hello!&num=42 HTTP/1.1
Host: example.com
User-Agent: Mozilla/5.0
Accept: text/plain
Accept-Language: en-us
Connection: keep-alive
\end{lstlisting}
\end{affianco}

Se l'utente scrive \code{Hello!} e \code{42}, la URL della richiesta diventa:

\begin{center}
  \code{/handler.html?msg=Hello!\&num=42}
\end{center}

La parte dopo il punto interrogativo si chiama \term{query string}, e
\code{msg} e \code{num} si dicono \term{parametri di query}.
\end{esempio}

\begin{figure}[H]
  \centering
  \begin{tikzpicture}[
      part/.style={font=\FiraCode\large, inner sep=1pt, outer sep=0pt},
      lab/.style={font=\Lato\scriptsize, align=center, inner sep=2pt}]
    \node[part, text=codekey]                       (sch) {https://};
    \node[part, text=codetag,    right=0pt of sch]  (dom) {app.com};
    \node[part, text=codestring, right=0pt of dom]  (pth) {/foo/bar.html};
    \node[part, text=codeattr,   right=0pt of pth]  (qry) {?key1=val1};
    \node[part, text=codenum,    right=0pt of qry]  (amp) {\&};
    \node[part, text=codeattr,   right=0pt of amp]  (qr2) {key2=val2};

    \node[lab, text=codekey,  above=30pt of sch] (l1) {Schema};
    \node[lab, text=codeattr, above=17pt of qry] (l4) {Parametri di query};
    \node[lab, text=codenum,  above=30pt of amp] (l5) {Separatore};
    \draw[wtdash] (l1.south) -- (sch.north);
    \draw[wtdash] (l4.south) -- (qry.north);
    \draw[wtdash] (l5.south) -- (amp.north);

    \node[lab, text=codetag,    below=17pt of dom] (l2) {Nome di dominio};
    \node[lab, text=codestring, below=17pt of pth] (l3) {Percorso};
    \draw[wtdash] (l2.north) -- (dom.south);
    \draw[wtdash] (l3.north) -- (pth.south);
  \end{tikzpicture}
  \caption{Una URL con parametri di query.}
  \label{fig:url-query}
\end{figure}

I parametri di query sono informazioni aggiuntive fornite al server, sotto forma
di coppie chiave/valore separate da \code{\&}. Il server può usarli per svolgere
elaborazioni prima di restituire la risorsa. Un esempio reale è la URL
\urlx{https://motore-di-ricerca.com/search?q=web+technologies&lang=it}.

\subsection{Invio con POST}

Se il metodo è \code{POST}, gli input vengono inviati nel \term{corpo} della
richiesta, e la URL della richiesta resta semplicemente \code{/handler.html}.

\begin{esempio}{La stessa richiesta, con il metodo POST}
\begin{affianco}[0.53]
\begin{lstlisting}[style=html, name=form-post.html]
<form action="/handler.html" method="POST">
  Messaggio: <input type="text" name="msg"><br>
  Numero: <input type="number" name="num"><br>
  <input type="submit">
</form>
\end{lstlisting}
\begin{browser}[form-post.html]
  Messaggio: \wtcampo[2.6cm]{Hello!}\\[3pt]
  Numero: \wtcampo[2.6cm]{42}\\[5pt]
  \wtbottone{Invia}
\end{browser}
\risultato
\begin{lstlisting}[style=http, name=Richiesta generata]
POST /handler.html HTTP/1.1
Host: example.com
User-Agent: Mozilla/5.0
Accept: text/plain
Accept-Language: en-us
Connection: keep-alive

msg=Hello!&num=42
\end{lstlisting}
\wtnota{I dati non sono più nella URL, ma nel \textbf{corpo},\\dopo la riga
        vuota che chiude le intestazioni.}
\end{affianco}
\end{esempio}

\warning{Meglio GET o POST?}{
  Con \code{GET} i dati finiscono nella URL: restano nella cronologia del
  browser, nei segnalibri, nei log del server e nell'header \code{Referer}
  inviato ai siti esterni. Non va quindi \textbf{mai} usato \code{GET} per
  password o dati personali. \code{GET} è indicato per ricerche e filtri, cioè
  quando avere una URL condivisibile che riproduce lo stesso risultato è un
  vantaggio; \code{POST} è indicato quando l'invio \textbf{modifica} qualcosa sul
  server.

  Attenzione però: \code{POST} non è \guillemotleft più sicuro\guillemotright\ in senso stretto. Il
  corpo della richiesta viaggia in chiaro esattamente come la URL. Ciò che
  protegge davvero i dati è HTTPS.}

\subsection{Codifica delle URL}

L'input di un form viene codificato come una stringa di coppie chiave-valore
separate da \code{\&}. Che cosa succede se un utente inserisce proprio dei
caratteri speciali, come \code{\&}?

Questi caratteri vengono sostituiti da terne del tipo \code{\%XX}, dove
\code{XX} sono due cifre esadecimali che rappresentano il carattere sostituito
nella tabella ASCII. Gli spazi possono essere sostituiti da \code{\%20} oppure
dal simbolo \code{+}.

\begin{esempio}{Percent encoding}
\begin{affianco}[0.5]
\begin{lstlisting}[style=html, name=encoding.html]
<form>
  <label for="msg">Messaggio:</label>
  <input id="msg" name="msg" type="text">
  <input type="submit">
</form>
\end{lstlisting}
\risultato
\begin{browser}[encoding.html]
  Messaggio: \wtcampo[2.8cm]{Tom \& Jerry} \wtbottone{Invia}
\end{browser}
\wtnota{Dato inviato: \code{msg=Tom+\%26+Jerry}}
\end{affianco}

Lo spazio è diventato \code{+} e la \code{\&} è diventata \code{\%26}, perché
\code{\&} vale \code{26} in esadecimale nella tabella ASCII.
\end{esempio}

\begin{center}
\renewcommand{\arraystretch}{1.2}
{\small
\begin{tabular}{@{}rrcl@{\hspace{1.2cm}}rrcl@{}}
  \textbf{Dec} & \textbf{Hex} & \textbf{Car.} & \textbf{Descrizione} &
  \textbf{Dec} & \textbf{Hex} & \textbf{Car.} & \textbf{Descrizione} \\
  32 & 20 &           & spazio              & 43 & 2B & \code{+} & più \\
  33 & 21 & \code{!}  & punto esclamativo   & 44 & 2C & \code{,} & virgola \\
  34 & 22 & \code{"}  & virgolette doppie   & 45 & 2D & \code{-} & meno \\
  35 & 23 & \code{\#} & cancelletto         & 46 & 2E & \code{.} & punto \\
  36 & 24 & \code{\$} & dollaro             & 47 & 2F & \code{/} & barra \\
  37 & 25 & \code{\%} & percento            & 91 & 5B & \code{[} & quadra aperta \\
  38 & 26 & \code{\&} & e commerciale       & 92 & 5C & \code{\textbackslash} & barra rovesciata \\
  39 & 27 & \code{'}  & apostrofo           & 93 & 5D & \code{]} & quadra chiusa \\
  40 & 28 & \code{(}  & tonda aperta        & 94 & 5E & \code{\textasciicircum} & circonflesso \\
  41 & 29 & \code{)}  & tonda chiusa        & 95 & 5F & \code{\_} & trattino basso \\
  42 & 2A & \code{*}  & asterisco           & 96 & 60 & \code{`} & accento grave \\
\end{tabular}}
\end{center}

\subsection{Caselle di spunta}

Le \term{checkbox} permettono di selezionare zero o più opzioni. Il
comportamento all'invio dipende da quante ne sono state spuntate.

\begin{esempio}{Comportamento delle checkbox all'invio}
\begin{affianco}[0.58]
\codicepiccolo
\begin{lstlisting}[style=html, name=checkbox.html]
<form>
  <p>Quali esami sosterrai?</p>
  <input type="checkbox" name="exams" value="web" id="web_tech">
  <label for="web_tech">Tecnologie Web</label><br>
  <input type="checkbox" name="exams" value="pl2" id="pl2">
  <label for="pl2">Linguaggi di Programmazione II</label><br><br>
  <input type="submit">
</form>
\end{lstlisting}
\risultato
\begin{browser}[checkbox.html]
  Quali esami sosterrai?\\[2pt]
  \wtcasella{1} Tecnologie Web\\
  \wtcasella{0} Linguaggi di Programmazione II\\[3pt]
  \wtbottone{Invia}
\end{browser}
\wtnota{All'invio: \code{exams=web}}
\begin{browser}[checkbox.html]
  Quali esami sosterrai?\\[2pt]
  \wtcasella{1} Tecnologie Web\\
  \wtcasella{1} Linguaggi di Programmazione II\\[3pt]
  \wtbottone{Invia}
\end{browser}
\wtnota{All'invio: \code{exams=web\&exams=pl2}}
\begin{browser}[checkbox.html]
  Quali esami sosterrai?\\[2pt]
  \wtcasella{0} Tecnologie Web\\
  \wtcasella{0} Linguaggi di Programmazione II\\[3pt]
  \wtbottone{Invia}
\end{browser}
\wtnota{All'invio: \code{""} (stringa vuota)}
\end{affianco}

Da notare: più caselle possono condividere lo \textbf{stesso} \attr{name}. In tal
caso il nome compare più volte nei dati inviati, una per ogni casella spuntata. E
se l'utente non spunta nulla, il parametro \textbf{non viene inviato affatto}: il
server deve essere preparato a non trovarlo.
\end{esempio}

\subsection{Pulsanti radio}

I \term{radio button} permettono di scegliere \textbf{una sola} fra più opzioni.
Le opzioni che condividono lo stesso \attr{name} formano un gruppo mutuamente
esclusivo.

\begin{esempio}{Un gruppo di radio button}
\begin{affianco}[0.58]
\codicepiccolo
\begin{lstlisting}[style=html, name=radio.html]
<form>
  <p>Qual e' il tuo corso preferito?</p>
  <input type="radio" name="fav" value="web" id="web_tech">
  <label for="web_tech">Tecnologie Web</label><br>
  <input type="radio" name="fav" value="net" id="net">
  <label for="net">Reti di Calcolatori</label><br>
  <input type="radio" name="fav" value="se" id="se">
  <label for="se">Ingegneria del Software</label><br><br>
  <input type="submit" value="Invia la tua opinione">
</form>
\end{lstlisting}
\risultato
\begin{browser}[radio.html]
  Qual è il tuo corso preferito?\\[2pt]
  \wtradio{1} Tecnologie Web\\
  \wtradio{0} Reti di Calcolatori\\
  \wtradio{0} Ingegneria del Software\\[4pt]
  \wtbottone{Invia la tua opinione}
\end{browser}
\wtnota{All'invio: \code{fav=web}}
\end{affianco}

Selezionando la prima opzione, all'invio si ottiene \code{fav=web}. Selezionarne
un'altra \textbf{deseleziona} automaticamente la precedente, perché le tre
opzioni condividono lo stesso \attr{name}.

Si noti anche che l'attributo \attr{value} su un \code{input type="submit"}
cambia l'\textbf{etichetta del pulsante}, mentre sugli altri input rappresenta il
valore inviato.
\end{esempio}

\subsection{Date e orari}

Esistono tipi di input dedicati alle date e agli orari.

\begin{affianco}[0.5]
\begin{lstlisting}[style=html, name=date.html]
<form>
  <input type="date"><br>
  <input type="week"><br>
  <input type="month"><br>
  <input type="time"><br>
  <input type="datetime-local"><br>
</form>
\end{lstlisting}
\risultato
\begin{browser}[date.html]
  \wtcampo[3.4cm]{17/09/2026 \hfill\tiny$\square$}\\[2pt]
  \wtcampo[3.4cm]{Settimana 38, 2026 \hfill\tiny$\square$}\\[2pt]
  \wtcampo[3.4cm]{settembre 2026 \hfill\tiny$\square$}\\[2pt]
  \wtcampo[3.4cm]{11:42 \hfill\tiny$\circ$}\\[2pt]
  \wtcampo[3.4cm]{17/09/2026, 09:45 \hfill\tiny$\square$}
\end{browser}
\wtnota{Valori già inseriti dall'utente; l'aspetto varia da browser a browser.}
\end{affianco}

\info{Il browser fa il lavoro sporco}{
  Usando \code{type="date"} il browser fornisce gratuitamente un calendario a
  comparsa, la validazione del formato e la localizzazione (un utente italiano
  vedrà \code{gg/mm/aaaa}, uno statunitense \code{mm/dd/yyyy}), mentre al server
  arriva sempre il formato normalizzato \code{aaaa-mm-gg}. Scrivere a mano lo
  stesso comportamento con un \code{type="text"} richiederebbe molto codice, e
  funzionerebbe peggio.}

\subsection{Menu a tendina}

L'elemento \tg{select} definisce un \term{menu a tendina}. L'attributo
\attr{multiple} permette di selezionare più di un'opzione.

\begin{esempio}{Un menu a tendina}
\begin{affianco}[0.55]
\begin{lstlisting}[style=html, name=select.html]
<form>
  <label for="dest">Seleziona la destinazione:</label>
  <select name="destination" id="dest">
    <option value="Rome">Roma</option>
    <option value="Naples">Napoli</option>
    <option value="Paris">Parigi</option>
  </select><br><br>
  <input type="submit">
</form>
\end{lstlisting}

\risultato
\begin{browser}[select.html]
  Seleziona la destinazione:\\[2pt]
  \wttendina[1.8cm]{Roma}\\[-3.2pt]
  {\setlength{\fboxsep}{0.6pt}\fcolorbox{neutral!55!white}{white}{%
    \parbox{\dimexpr1.8cm+6pt\relax}{\setlength{\fboxsep}{1.5pt}%
      \colorbox{white}{\makebox[\dimexpr1.8cm+3pt\relax][l]{\rule{0pt}{8pt}\wtuifont Roma}}\\[-1pt]
      \colorbox{dtblue}{\makebox[\dimexpr1.8cm+3pt\relax][l]{\rule{0pt}{8pt}\wtuifont\color{white}Napoli}}\\[-1pt]
      \colorbox{white}{\makebox[\dimexpr1.8cm+3pt\relax][l]{\rule{0pt}{8pt}\wtuifont Parigi}}}}}\\[5pt]
  \wtbottone{Invia}
\end{browser}
\wtnota{Il menu aperto: il mouse è su \textit{Napoli}.}
\end{affianco}

Da notare la differenza fra il testo mostrato all'utente (\textit{Roma}) e il
valore inviato al server (\code{Rome}), specificato tramite \attr{value}.
\end{esempio}

\subsection{E molto altro}

Esistono numerosi altri tipi di input: selettore di colore, selettore di file,
liste di suggerimenti (\code{datalist}), campi per e-mail, URL e numeri di
telefono, cursori a scorrimento.

\begin{affianco}[0.55]
\begin{lstlisting}[style=html, name=colore.html]
<form method="POST">
  <label for="col">Scegli il tuo colore preferito:</label>
  <input name="color" id="col" type="color"><br><br>
  <input type="submit">
</form>
\end{lstlisting}
\risultato
\begin{browser}[colore.html]
  Scegli il tuo colore preferito:
  \fcolorbox{neutral!55!white}{neutral!8!white}{%
    \textcolor[HTML]{1D8FA8}{\rule[-1.5pt]{0.9cm}{8pt}}}\\[8pt]
  \wtbottone{Invia}
\end{browser}
\wtnota{Cliccando sul campione si apre il selettore di colore del sistema;\\
        il valore inviato è in esadecimale, per esempio \code{color=\%231d8fa8}.}
\end{affianco}

Per un riferimento completo conviene consultare la documentazione MDN.

\subsection{Raggruppare gli input}

Gli elementi \tg{fieldset} e \tg{legend} permettono di raggruppare logicamente
gli input di un form e di dare un titolo a ciascun gruppo.

\begin{esempio}{Un form suddiviso in sezioni}
\begin{affianco}[0.55]
\begin{lstlisting}[style=html, name=fieldset.html]
<form>
  <fieldset>
    <legend>Dati anagrafici:</legend>
    <label for="fname">Nome:</label><br>
    <input type="text" id="fname" name="fname"><br>
    <label for="lname">Cognome:</label><br>
    <input type="text" id="lname" name="lname">
  </fieldset>
  <fieldset>
    <legend>Registrazione esame:</legend>
    <label for="grade">Voto:</label><br>
    <input type="number" id="grade" name="grade"><br>
    <label for="date">Data:</label><br>
    <input type="date" id="date" name="date">
  </fieldset><br>
  <input type="submit" value="Invia">
</form>
\end{lstlisting}
\risultato
\begin{browser}[fieldset.html]
  \begin{wtfieldset}{Dati anagrafici:}
    Nome:\\ \wtcampo[3cm]{John}\\
    Cognome:\\ \wtcampo[3cm]{Doe}
  \end{wtfieldset}
  \begin{wtfieldset}{Registrazione esame:}
    Voto:\\ \wtcampo[3cm]{30}\\
    Data:\\ \wtcampo[3cm]{30/09/2024}
  \end{wtfieldset}
  \vspace{4pt}\wtbottone{Invia}
\end{browser}
\wtnota{All'invio: \code{fname=John\&lname=Doe\&grade=30\&date=2024-09-30}}
\end{affianco}

Il \tg{fieldset} raggruppa gli input \textit{visivamente} e
\textit{semanticamente}, ma non ha alcun effetto sui dati inviati: questi
restano una lista piatta di coppie nome/valore.
\end{esempio}

\section{Organizzare il contenuto}

Il contenuto di una pagina web può essere organizzato (raggruppato) in parti.
Questo si può fare in due modi:

\begin{itemize}
  \item con le \term{divisioni} \tg{div};
  \item con i \term{tag semantici} come \tg{header}, \tg{nav}, \tg{main},
        \tg{article}, \tg{section}, \tg{aside} e \tg{footer}.
\end{itemize}

\subsection{Divisioni}

Le divisioni \tg{div} erano il principale strumento di raggruppamento nelle
versioni più vecchie di HTML, prima dell'introduzione dei tag semantici. Non
hanno alcuna semantica particolare, se non quella di raggruppare contenuti in
qualche modo correlati fra loro.

\begin{affianco}[0.5]
\begin{lstlisting}[style=html, name=divisioni.html]
<div>
  <h1>Sezione 1</h1>
  <p>Contenuto della Sezione 1</p>
</div>
<div>
  <h1>Sezione 2</h1>
  <p>Contenuto della Sezione 2</p>
</div>
\end{lstlisting}
\risultato
\begin{browser}[divisioni.html]
  {\Lora\bfseries\LARGE Sezione 1}\\[4pt]
  Contenuto della Sezione 1\\[8pt]
  {\Lora\bfseries\LARGE Sezione 2}\\[4pt]
  Contenuto della Sezione 2
\end{browser}
\wtnota{Le \code{<div>} non si vedono: raggruppano senza cambiare l'aspetto.}
\end{affianco}

\subsection{Lo span}

La \tg{div} ha un equivalente per il testo: l'elemento \term{span}. Anche
\tg{span} è un contenitore \textbf{generico e privo di significato}, con una
sola differenza, ma decisiva: \tg{div} è un elemento di tipo \textbf{blocco}
(occupa una riga tutta sua), mentre \tg{span} è \textbf{inline}, cioè resta nel
flusso del testo come una parola qualsiasi.

Serve quando si vuole identificare una \textbf{porzione di testo} all'interno di
un elemento più grande, per darle uno stile o per manipolarla con JavaScript,
e non esiste alcun tag che ne descriva il significato.

\begin{affianco}[0.55]
\begin{lstlisting}[style=html, name=span.html]
<p>
  Il totale è <span class="prezzo">19,99 €</span>,
  spedizione <span class="gratis">inclusa</span>.
</p>
<p>
  Un <span>span senza classe</span> non cambia nulla.
</p>
\end{lstlisting}
\risultato
\begin{browser}[span.html]
  Il totale è {\bfseries\color{css-red}19,99 €}, spedizione
  {\color{css-green}inclusa}.\\[6pt]
  Un span senza classe non cambia nulla.
\end{browser}
\wtnota{I primi due \tg{span} sono colorati dal CSS (che vedremo nel prossimo
        capitolo); il terzo, senza stile, è invisibile.}
\end{affianco}

Di per sé \tg{span} non produce \textbf{alcun} effetto visibile: è un gancio a
cui appendere una classe o un identificatore.

\warning{Prima cerca un tag che dica qualcosa}{
  \tg{span} e \tg{div} sono privi di semantica: per un lettore di schermo o per
  un motore di ricerca non significano nulla. Vanno quindi usati soltanto
  quando nessun altro elemento è adatto. Se il testo va \textbf{enfatizzato}
  esiste \tg{em}, se è \textbf{importante} esiste \tg{strong}, per una sigla
  esiste \tg{abbr}, per una data \tg{time}, per il codice \tg{code}. La regola
  è sempre la stessa: prima il significato, poi l'aspetto.}

\subsection{Tag semantici}

I \term{tag semantici} descrivono il significato del loro contenuto ai browser,
agli sviluppatori e ai programmi che analizzano la pagina.

Le slide del corso li riassumono così: \tg{nav} contiene i collegamenti di
navigazione, \tg{main} il contenuto principale, \tg{article} contenuto
indipendente e autocontenuto, \tg{aside} contenuto solo tangenzialmente
correlato, mentre \tg{header}, \tg{footer} e \tg{section} «si spiegano da soli».
Vale la pena, invece, spiegarli tutti con un po' più di precisione.

\begin{center}
\renewcommand{\arraystretch}{1.45}
\begin{tabular}{@{}l@{\hspace{0.8cm}}>{\raggedright\arraybackslash}p{7.3cm}@{\hspace{0.6cm}}>{\raggedright\arraybackslash}p{6.2cm}@{}}
  \textbf{Tag} & \textbf{Che cosa rappresenta} & \textbf{Contenuto tipico} \\
  \tg{header} & L'\textbf{intestazione} introduttiva della pagina, oppure di un
                \tg{article} o di una \tg{section}. &
                Logo, titolo del sito, titolo e autore di un articolo, a volte
                il menu. \\
  \tg{nav}    & Un blocco di \textbf{collegamenti di navigazione} principali. &
                Menu del sito, indice di una pagina lunga, briciole di pane.
                Non ogni gruppo di link: solo quelli per muoversi nel sito. \\
  \tg{main}   & Il \textbf{contenuto principale} e specifico della pagina.
                Ne esiste \textbf{uno solo} per pagina. &
                Tutto ciò che non si ripete uguale su ogni pagina: esclusi
                intestazione, menu e piè di pagina del sito. \\
  \tg{article}& Un contenuto \textbf{autocontenuto}, che avrebbe senso anche
                preso da solo, fuori dalla pagina. &
                Un post di un blog, una notizia, un commento, la scheda di un
                prodotto. \\
  \tg{section}& Un \textbf{raggruppamento tematico} di contenuti, parte di un
                insieme più grande; di norma ha un proprio titolo. &
                I capitoli di un documento, le sezioni «Pro» e «Contro» di un
                articolo. \\
  \tg{aside}  & Contenuto \textbf{correlato solo di lato} al contenuto
                principale: toglierlo non ne cambia il senso. &
                Barra laterale, box di approfondimento, glossario, pubblicità. \\
  \tg{footer} & Il \textbf{piè di pagina} della pagina, o di un \tg{article}
                o di una \tg{section}. &
                Copyright, autore, contatti, link a informazioni legali. \\
\end{tabular}
\end{center}

Tre osservazioni aiutano a non confonderli.

\begin{itemize}
  \item \tg{header} e \tg{footer} \textbf{non} sono per forza unici: la pagina ha
        la sua intestazione, ma ogni \tg{article} può avere la propria. Il nome
        non indica una posizione sullo schermo, ma un \textbf{ruolo} rispetto al
        contenuto che li contiene.
  \item La differenza fra \tg{article} e \tg{section} si decide con una domanda:
        \textit{questo pezzo avrebbe senso pubblicato da solo, altrove?} Se sì è
        un \tg{article}; se è solo una parte di un discorso più ampio è una
        \tg{section}. Un articolo può contenere sezioni, e (più raramente) una
        sezione può contenere articoli, come l'elenco delle notizie di una
        homepage.
  \item Se serve un contenitore soltanto per applicare uno stile, senza alcun
        significato particolare, la scelta giusta resta \tg{div}: una
        \tg{section} senza titolo è quasi sempre una \tg{div} travestita.
\end{itemize}

\begin{figure}[H]
  \centering
  \begin{tikzpicture}[font=\FiraCode\footnotesize,
      b/.style={rounded corners=2pt, line width=0.8pt, anchor=north west}]
    \node[b, draw=primary!70!black, fill=primary!4!white,
          minimum width=10.4cm, minimum height=6.7cm] (body) at (0,0) {};
    \node[anchor=north west, text=primary!70!black] at (0.1,-0.08) {<body>};

    \node[b, draw=accent!60!black, fill=accent!8!white,
          minimum width=9.6cm, minimum height=1.0cm] at (0.4,-0.6) {};
    \node[anchor=north west, text=accent!60!black] at (0.5,-0.68) {<header>};

    \node[b, draw=orange-gradient-6, fill=orange-gradient-1!35!white,
          minimum width=9.6cm, minimum height=0.85cm] at (0.4,-1.7) {};
    \node[anchor=north west, text=orange-gradient-6!70!black] at (0.5,-1.78) {<nav>};

    \node[b, draw=primary!60!black, fill=primary!8!white,
          minimum width=9.6cm, minimum height=3.05cm] at (0.4,-2.65) {};
    \node[anchor=north west, text=primary!60!black] at (0.5,-2.73) {<main>};

    \node[b, draw=green-gradient-6, fill=green-gradient-1!45!white,
          minimum width=6.2cm, minimum height=2.15cm] at (0.8,-3.3) {};
    \node[anchor=north west, text=green-gradient-6!55!black] at (0.9,-3.38) {<article>};

    \node[b, draw=green-gradient-6, fill=white,
          minimum width=5.6cm, minimum height=0.6cm] at (1.1,-3.9) {};
    \node[anchor=north west, text=green-gradient-6!55!black] at (1.2,-3.96) {<section>};
    \node[b, draw=green-gradient-6, fill=white,
          minimum width=5.6cm, minimum height=0.6cm] at (1.1,-4.6) {};
    \node[anchor=north west, text=green-gradient-6!55!black] at (1.2,-4.66) {<section>};

    \node[b, draw=azure-gradient-6, fill=azure-gradient-1!45!white,
          minimum width=2.7cm, minimum height=2.15cm] at (7.3,-3.3) {};
    \node[anchor=north west, text=azure-gradient-6!55!black] at (7.4,-3.38) {<aside>};

    \node[b, draw=accent!60!black, fill=accent!8!white,
          minimum width=9.6cm, minimum height=0.85cm] at (0.4,-5.85) {};
    \node[anchor=north west, text=accent!60!black] at (0.5,-5.93) {<footer>};
  \end{tikzpicture}
  \caption{Un tipico impianto di pagina realizzato con tag semantici.}
  \label{fig:semantic}
\end{figure}

\begin{esempio}{Una pagina strutturata semanticamente}
\begin{affianco}[0.55]
\codicepiccolo
\begin{lstlisting}[style=html, name=semantica-completa.html]
<header>
  <h1>Tecnologie Web!</h1>
  <p>Questa pagina contiene contenuti sulle Tecnologie Web.</p>
</header>
<main>
  <article>
    <h2>Gli elementi semantici sono una buona cosa</h2>
    <p>Parliamo degli elementi semantici.</p>
    <section>
      <h3>Pro</h3>
      <p>Veicolano piu' informazione.</p>
    </section>
    <section>
      <h3>Contro</h3>
      <p>Letteralmente nessuno.</p>
    </section>
  </article>
  <article>
    <h2>HTML e' bello</h2>
    <p>E non hai ancora visto lo stile e gli script!</p>
  </article>
</main>
<footer> &copy; Corso di Tecnologie Web, 2024. </footer>
\end{lstlisting}
\risultato
\begin{browser}[semantica-completa.html]
  {\Lora\bfseries\LARGE Tecnologie Web!}\\[3pt]
  {\small Questa pagina contiene contenuti sulle Tecnologie Web.}\\[6pt]
  {\Lora\bfseries\Large Gli elementi semantici sono una buona cosa}\\[3pt]
  {\small Parliamo degli elementi semantici.}\\[5pt]
  {\Lora\bfseries\large Pro}\\[2pt]
  {\small Veicolano più informazione.}\\[5pt]
  {\Lora\bfseries\large Contro}\\[2pt]
  {\small Letteralmente nessuno.}\\[6pt]
  {\Lora\bfseries\Large HTML è bello}\\[3pt]
  {\small E non hai ancora visto lo stile e gli script!}\\[6pt]
  {\small \copyright\ Corso di Tecnologie Web, 2024.}
\end{browser}
\end{affianco}

Nel risultato non c'è traccia di \tg{header}, \tg{article} o \tg{footer}: senza
CSS i tag semantici non cambiano l'aspetto della pagina (si vedono solo i titoli
e i paragrafi). Il loro valore sta tutto nell'informazione che aggiungono, come
spiega il riquadro seguente.
\end{esempio}

\info{Perché preferire i tag semantici alle divisioni}{
  Visivamente \tg{div} e \tg{section} sono identici: nessuno dei due, di per sé,
  cambia l'aspetto della pagina. La differenza sta nell'\textbf{informazione}. Un
  lettore di schermo può offrire il comando \guillemotleft vai al contenuto principale\guillemotright\
  solo se esiste un \tg{main}; un motore di ricerca capisce che cosa indicizzare
  come articolo solo se esiste un \tg{article}; la modalità lettura del browser
  funziona per la stessa ragione. Una pagina fatta solo di \tg{div} funziona, ma
  è muta.}

\section{Il documento come albero}

Gli elementi HTML sono \textbf{annidati} l'uno dentro l'altro: un \tg{body}
contiene un \tg{main}, che contiene un \tg{article}, che contiene un \tg{h1} e un
\tg{p}. Questa struttura di contenimento è, formalmente, un \term{albero}.

\begin{figure}[H]
  \centering
  \begin{minipage}[c]{0.40\linewidth}
\begin{lstlisting}[style=html, basicstyle=\linespread{1.1}\scriptsize\FiraCode]
<!DOCTYPE html>
<html lang="it">
  <head>
    <title>Il Libro</title>
    <meta charset="UTF-8">
  </head>
  <body>
    <main>
      <article>
        <h1>Titolo</h1>
        <p>Corpo</p>
      </article>
    </main>
    <footer>
      <p>&copy; TecWeb</p>
    </footer>
  </body>
</html>
\end{lstlisting}
  \end{minipage}\hfill
  \begin{minipage}[c]{0.57\linewidth}
    \centering
    \begin{tikzpicture}[font=\FiraCode\scriptsize,
        level distance=0.9cm,
        n/.style={rounded corners=2pt, draw=primary!60!black, line width=0.7pt,
                  fill=primary!8!white, inner sep=2.5pt},
        edge from parent/.style={draw=neutral!50!black, line width=0.6pt,
                                 edge from parent path=
                                 {(\tikzparentnode.south) -- (\tikzchildnode.north)}},
        level 1/.style={sibling distance=3.5cm},
        level 2/.style={sibling distance=1.8cm},
        level 3/.style={sibling distance=1.25cm}]
      \node[n] {<html>}
        child { node[n] {<head>}
          child { node[n] {<title>} }
          child { node[n] {<meta>} } }
        child { node[n] {<body>}
          child { node[n] {<main>}
            child { node[n] {<article>}
              child { node[n] {<h1>} }
              child { node[n] {<p>} } } }
          child { node[n] {<footer>}
            child { node[n] {<p>} } } };
    \end{tikzpicture}
  \end{minipage}
  \caption{Lo stesso documento visto come sorgente e come albero.}
  \label{fig:dom-tree}
\end{figure}

Questa non è una curiosità teorica: è il modello su cui si fonda tutto il resto
del corso. Il browser, quando analizza il documento, costruisce esattamente
questo albero in memoria, e lo chiama \term{DOM} (\textit{Document Object
Model}). I selettori CSS servono a individuare nodi in questo albero, e
JavaScript lo modifica per rendere la pagina dinamica.

\error{Annidamento sbagliato}{
  Se si scrive \code{<p>testo <strong>importante</p></strong>}, i tag si
  \guillemotleft incrociano\guillemotright\ e la struttura non è più un albero. I browser non segnalano
  l'errore: cercano di indovinare che cosa intendevamo, e ognuno può indovinare
  in modo diverso. Il risultato è una pagina che appare corretta su un browser e
  sbagliata su un altro. I tag vanno sempre chiusi nell'ordine inverso rispetto a
  quello di apertura.}

\section{Gli strumenti per sviluppatori del browser}

I browser moderni includono numerose funzionalità pensate per chi sviluppa. Si
accede a questi \term{DevTools} premendo il tasto \code{F12}.

Le funzionalità principali sono:

\begin{itemize}
  \item \textbf{ispezione} dei documenti HTML: si vede l'albero del DOM così come
        il browser lo ha effettivamente costruito, e lo si può modificare dal
        vivo;
  \item \textbf{analisi di rete}: il dettaglio di tutte le richieste e risposte
        HTTP coinvolte, con header, corpo, codice di stato e tempi;
  \item \textbf{profilazione}: misura delle prestazioni e dei tempi di
        caricamento;
  \item \textbf{debug} sia della presentazione grafica sia degli script.
\end{itemize}

\subsection{Ispezionare la pagina}

La scheda \textit{Analisi pagina} (in Chrome \textit{Elements}) mostra l'albero
del DOM costruito dal browser. Passando con il mouse su un nodo, l'elemento
corrispondente si evidenzia nella pagina; nella parte bassa del pannello si
vedono gli stili che gli vengono applicati, che impareremo a leggere nel
prossimo capitolo. Il DOM mostrato qui è quello \textbf{effettivo}: se il
browser ha dovuto correggere un annidamento sbagliato, l'albero mostra la sua
correzione, non il sorgente originale.

\subsection{Analisi di rete}

La scheda \textit{Rete} (in Chrome \textit{Network}) registra \textbf{tutte le
richieste HTTP} effettuate dalla pagina. La figura riproduce ciò che mostra
Firefox caricando la pagina \code{image.html} dell'esempio sulle immagini.

\begin{figure}[H]
  \centering
  \begin{minipage}[t]{3.6cm}\vspace{0pt}
    \begin{browser}[Fu-Tzu]
      {\Lora\bfseries\large Fu-Tzu}\\[3pt]
      {\scriptsize Un ritratto del maestro, generato da DALL-E, è disponibile
      qui di seguito:}\\[3pt]
      \includegraphics[width=2.4cm]{sources/fu-tzu.jpg}
    \end{browser}
  \end{minipage}\hfill
  \begin{minipage}[t]{15.4cm}\vspace{0pt}
  \begin{devtools}{Rete}
    {\Lato\scriptsize\setlength{\tabcolsep}{5pt}%
    \colorbox{dtbar}{\makebox[\dimexpr\linewidth-2\fboxsep][l]{%
      \color{neutral!45!black}\quad Filtra URL\hfill
      \color{neutral!20!black}\wtcasella{1}\ Disattiva cache\quad Nessun limite\quad}}\par
    \noindent\hspace*{6pt}{\color{neutral!30!black}Tutti\quad
      \colorbox{dtline}{HTML}\quad CSS\quad JS\quad XHR\quad Caratteri\quad
      \colorbox{dtline}{Immagini}\quad Media\quad WS\quad Altro}\par\vspace{2pt}
    {\color{dtline}\hrule}
    \noindent\begin{tabular*}{\linewidth}{@{\hspace{8pt}}l@{\extracolsep{\fill}}llllll@{\hspace{8pt}}}
      \rule{0pt}{10pt}\textbf{Stato} & \textbf{Metodo} & \textbf{Dominio} &
      \textbf{File} & \textbf{Iniziatore} & \textbf{Tipo} & \textbf{Dimensione} \\ \hline
      \rule{0pt}{10pt}\dtstato{200} & GET & 127.0.0.1:3000 & image.html & document & html & 373 B \\
      \dtstato{200} & GET & 127.0.0.1:3000 & pic.jpg & img & jpeg & 109,90 kB \\
      \dtstato{200} & GET & 127.0.0.1:3000 & favicon.ico &
        \textcolor{dtblue}{\uline{FaviconLoader.sys.mjs}} & vnd.microsoft.icon & 15,41 kB \\
    \end{tabular*}\par\vspace{10pt}
    {\color{dtline}\hrule}\vspace{2pt}
    \noindent\hspace*{6pt}{\color{neutral!25!black}3 richieste\quad$\vert$\quad
      125,68 kB di 126,25 kB trasferiti\quad$\vert$\quad Completato: 109 ms\quad$\vert$\quad
      \textcolor{dtblue}{DOMContentLoaded: 57 ms}\quad$\vert$\quad
      \textcolor{dtsel}{load: 61 ms}}\par\vspace{3pt}}
  \end{devtools}
  \end{minipage}
  \caption{La scheda \textit{Rete} di Firefox durante il caricamento di
           \code{image.html}: tre richieste per una sola pagina.}
  \label{fig:devtools-rete}
\end{figure}

Ogni riga è una richiesta, e le colonne raccontano che cosa è successo:

\begin{itemize}
  \item \textbf{Stato} e \textbf{Metodo}: il codice di stato della risposta e il
        verbo HTTP usato (qui sempre \code{GET}, tutto a buon fine con
        \code{200});
  \item \textbf{Dominio} e \textbf{File}: a chi è stata fatta la richiesta e
        quale risorsa è stata chiesta;
  \item \textbf{Iniziatore}: \textbf{chi} ha causato la richiesta. Il documento
        è stato richiesto dalla navigazione (\code{document}), l'immagine dal tag
        \tg{img} (\code{img}); la favicon, invece, l'ha chiesta il
        \textbf{browser di sua iniziativa} (\code{FaviconLoader}), pur non
        essendo dichiarata da nessuna parte nella pagina;
  \item \textbf{Tipo} e \textbf{Dimensione}: il tipo di contenuto ricevuto (dal
        \code{Content-Type} della risposta) e il suo peso.
\end{itemize}

La barra in fondo riassume il caricamento: numero di richieste, byte trasferiti
e tre tempi. \code{DOMContentLoaded} indica quando il browser ha finito di
analizzare l'HTML; \code{load} quando sono arrivate anche tutte le risorse
esterne, immagini comprese; \textit{Completato} quando è terminata l'ultima
richiesta. Facendo clic su una riga si aprono i dettagli: intestazioni di
richiesta e risposta, corpo, tempi di ciascuna fase.

La fila di pulsanti \textit{Tutti, HTML, CSS, JS, Immagini, \ldots} filtra le
richieste per tipo: su un sito reale, con decine di richieste, è indispensabile.

\subsubsection{La cache del browser}

Nella barra degli strumenti della figura c'è una casella spuntata:
\textit{Disattiva cache}. Per capire a che cosa serve bisogna sapere che il
browser, dopo aver scaricato una risorsa, ne conserva una copia in una memoria
locale, la \term{cache}, per non doverla scaricare di nuovo alla visita
successiva.

\dfn{Cache HTTP}{
  La \term{cache HTTP} è la copia locale delle risposte già ricevute che il
  browser conserva per riusarle. Quanto a lungo una risposta possa essere
  riusata, e a quali condizioni, lo decide il server tramite le intestazioni
  della risposta (per esempio \code{Cache-Control}, \code{ETag},
  \code{Last-Modified}).}

Ricaricando la pagina con la cache attiva, i casi tipici sono tre.

\begin{itemize}
  \item La risorsa è ancora \textbf{fresca}: il browser la prende dalla cache
        \textbf{senza fare alcuna richiesta}. Nella scheda \textit{Rete} la riga
        compare comunque, ma come trasferito è indicata la cache (Chrome scrive
        \textit{memory cache} o \textit{disk cache}) e il tempo è quasi nullo.
  \item La risorsa va \textbf{riconvalidata}: il browser chiede al server
        \guillemotleft è cambiata dall'ultima volta?\guillemotright\ con una
        \term{richiesta condizionale}. Se non è cambiata, il server risponde
        \code{304 Not Modified} con un corpo \textbf{vuoto}, e il browser usa la
        copia che ha già: viaggiano poche centinaia di byte invece dell'intero
        file.
  \item La risorsa è cambiata (o non è in cache): risposta \code{200} con il
        contenuto completo, come al primo caricamento.
\end{itemize}

\begin{affianco}[0.5]
\begin{lstlisting}[style=http, name=richiesta condizionale]
GET /02-HTML/pic.jpg HTTP/1.1
Host: 127.0.0.1:3000
If-None-Match: "1ad3e-18a9c2f0b40"
If-Modified-Since: Tue, 15 Sep 2026 10:12:08 GMT
\end{lstlisting}
\risultato
\begin{lstlisting}[style=http, name=risposta del server]
HTTP/1.1 304 Not Modified
ETag: "1ad3e-18a9c2f0b40"
Last-Modified: Tue, 15 Sep 2026 10:12:08 GMT
\end{lstlisting}
\wtnota{Nessun corpo: l'immagine non viene trasmessa di nuovo.}
\end{affianco}

Il meccanismo si regge su due coppie di intestazioni. Alla prima risposta il
server allega un'impronta della risorsa (\code{ETag}) e la data dell'ultima
modifica (\code{Last-Modified}); alla richiesta successiva il browser le rimanda
indietro in \code{If-None-Match} e \code{If-Modified-Since}, e il server
confronta.

\begin{figure}[H]
  \centering
  \begin{devtools}[width=15.4cm]{Rete}
    {\Lato\scriptsize\setlength{\tabcolsep}{6pt}%
    \colorbox{dtbar}{\makebox[\dimexpr\linewidth-2\fboxsep][l]{%
      \color{neutral!45!black}\quad Filtra URL\hfill
      \color{neutral!20!black}\wtcasella{0}\ Disattiva cache\quad Nessun limite\quad}}\par
    \noindent\begin{tabular*}{\linewidth}{@{\hspace{8pt}}l@{\extracolsep{\fill}}lllll@{\hspace{8pt}}}
      \rule{0pt}{10pt}\textbf{Stato} & \textbf{Metodo} & \textbf{File} &
      \textbf{Iniziatore} & \textbf{Trasferito} & \textbf{Dimensione} \\ \hline
      \rule{0pt}{10pt}\dtstato[dtwarn]{304} & GET & image.html & document & 212 B & 373 B \\
      \dtstato[dtwarn]{304} & GET & pic.jpg & img & 214 B & 109,90 kB \\
      \dtstato{200} & GET & favicon.ico & FaviconLoader.sys.mjs &
        \textit{in cache} & 15,41 kB \\
    \end{tabular*}\par\vspace{8pt}
    {\color{dtline}\hrule}\vspace{2pt}
    \noindent\hspace*{6pt}{\color{neutral!25!black}3 richieste\quad$\vert$\quad
      126,25 kB / 426 B trasferiti\quad$\vert$\quad Completato: 21 ms}\par\vspace{3pt}}
  \end{devtools}
  \caption{La stessa pagina ricaricata con la cache attiva (valori indicativi):
           documento e immagine vengono riconvalidati con \code{304}, la favicon
           non viene nemmeno richiesta.}
  \label{fig:devtools-cache}
\end{figure}

\warning{Perché durante lo sviluppo si disattiva la cache}{
  Mentre si scrive una pagina, la cache è più un ostacolo che un aiuto: si
  modifica il CSS, si ricarica, e il browser mostra ancora la versione vecchia
  presa dalla cache. Ci sono due rimedi.
  \begin{itemize}
    \item Spuntare \textit{Disattiva cache} nella scheda \textit{Rete}: finché i
          DevTools restano aperti, il browser ignora la cache e riscarica tutto.
    \item Eseguire una \textbf{ricarica forzata} con \code{Ctrl+Shift+R} (oppure
          \code{Ctrl+F5}; su macOS \code{Cmd+Shift+R}), che scavalca la cache
          per quel singolo caricamento.
  \end{itemize}
  In entrambi i casi nella scheda \textit{Rete} tutte le righe tornano a
  \code{200}, con i file trasferiti per intero.}

\info{I DevTools saranno il tuo migliore amico}{
  Chi sviluppa per il Web li tiene aperti quasi sempre. Vale la pena spendere
  mezz'ora a esplorarli con calma su un sito qualunque: cambiare un colore dal
  vivo nella scheda \textit{Inspector}, guardare l'ordine e i tempi delle
  richieste nella scheda \textit{Network}, provare a leggere una risposta HTTP
  reale con i suoi header. Molti problemi che sembrano misteriosi diventano
  evidenti dopo dieci secondi passati nel pannello giusto.}

\section{Riferimenti}

\begin{itemize}
  \item \textbf{Learn HTML} --- web.dev. Sezioni da 1 a 10 e da 12 a 14. \\
        \urlx{https://web.dev/learn/html}
  \item \textbf{Learn Forms} --- web.dev. Sezioni da 1 a 3. \\
        \urlx{https://web.dev/learn/forms}
  \item \textbf{The Input (Form Input) element} --- MDN web docs. \\
        \urlx{https://developer.mozilla.org/en-US/docs/Web/HTML/Element/input}
  \item \textbf{HTML Forms (overview)} --- W3Schools. \\
        \urlx{https://www.w3schools.com/html/html_forms.asp}
  \item \textbf{HTML Living Standard} --- WHATWG. \\
        \urlx{https://html.spec.whatwg.org/} \\
        Non è richiesto di imparare l'intera specifica: è sufficiente sapere che
        esiste e farsi un'idea di come è strutturata.
  \item \textbf{Browser DevTools} \\
        Firefox: \urlx{https://firefox-source-docs.mozilla.org/devtools-user/} \\
        Chrome: \urlx{https://developer.chrome.com/docs/devtools} \\
        Conviene prendere confidenza con i DevTools del browser che si preferisce.
\end{itemize}
